% ==============================================================================
% Fig.1–Fig.5 的图题与「注」——可直接粘贴进 main.tex
% ==============================================================================
% 规范：图内**不放任何**标题、结论句与脚注，只保留坐标轴、刻度数字、
%       panel 字母与图例；一切说明性文字写在下方这里。
% 图片文件：04_图/fig01_framework.png …… fig05_cost_waterfall.png
%          （矢量件在 04_图/pdf/，同名 .pdf 与 .svg）
% 建议宽度：Fig.1/2/4/5 用 \textwidth，Fig.3 用 \textwidth。
% ==============================================================================


% ---------------------------------------------------------------- 图 1
\begin{figure}[htbp]
  \centering
  \includegraphics[width=\textwidth]{fig01_framework.png}
  \caption{统一建模框架与四问递进关系}
  \label{fig:framework}
\end{figure}

\noindent 注：实心圆表示该问题具备此环节，空心圆表示尚未引入；圆点与条带颜色
沿用全文语义色板。左侧青色回路为\textbf{跨日滚动}（当日日末实际储电量回传次日
预测），右侧蓝色弧线为\textbf{日内滚动}（6:00 / 12:00 / 18:00 重解剩余时段）。
三点需注意：其一，问题一的附件 1 直接给定电价、负载与光伏预测功率，
\textbf{不含}预测、不确定性分析与风险修正环节；其二，问题二只有"预测层—计划层—
执行层"三层，当日计划一经冻结不再修改，\textbf{不含}滚动环节，滚动由问题三引入；
其三，问题三的 6 / 12 / 18 是\textbf{发布时刻}而非提前小时数——附件 3 每天在这四个
时刻各发布\textbf{整条 24 点功率曲线}，模型只正式提交当前 6 h 交付块，初始计划
$g^{0}$ 永久冻结。


% ---------------------------------------------------------------- 图 2
\begin{figure}[htbp]
  \centering
  \includegraphics[width=\textwidth]{fig02_p1_dispatch.png}
  \caption{问题一：低价购电—储能充电—高价放电的确定性调度机制}
  \label{fig:p1-dispatch}
\end{figure}

\noindent 注：(a)(b)(c)(d) 四幅共享同一时间轴，依次为电价、计划购电量、
充放电量与储电量，构成"价格信号 $\rightarrow$ 购电决策 $\rightarrow$ 储能动作
$\rightarrow$ 状态后果"的完整因果链。浅绿底纹为低谷窗口（电价 $\le 0.436$ 元/kWh，
即当日电价的 25\% 分位以下），浅赭底纹为高峰窗口（电价 $\ge 0.984$ 元/kWh，
即 75\% 分位以上）；两个阈值均由当日电价自身确定，未引入外部假设。
当日峰谷价差达 \textbf{3.8 倍}；低谷时段仅占全天 \textbf{25\%}，却承担了
\textbf{53\%} 的购电量；储电量在低谷抬升、高峰回落，两端分别触及 10800 kWh 与
1200 kWh 的运行边界。数据来源于附件 1 的确定性 MILP 求解结果，详见
支撑材料 \texttt{p1\_detail.csv}。


% ---------------------------------------------------------------- 图 3
\begin{figure}[htbp]
  \centering
  \includegraphics[width=\textwidth]{fig03_q80_mechanism.png}
  \caption{问题二：风险分位数取 $0.8$ 的经济学解释}
  \label{fig:q80}
\end{figure}

\noindent 注：正常购电价为 $p$、紧急购电价为 $5p$，两侧代价不对称。多买
1 kWh 的确定成本为 $p$，少买 1 kWh 的期望边际成本为 $5p\Pr(N>g)$；令二者相等，
由 $\dfrac{\partial}{\partial g}=p-5p\Pr(N>g)=0$ 得 $F_N(g)=\dfrac{5-1}{5}=0.8$，
即最优计划量应取净负荷的 80\% 分位数。(a) 为净负荷预测误差的经验分布，
红色部分为右侧 20\% 尾部（即触发 5 倍价紧急购电的区间），Q0.8 位置已标出；
(b) 为两条边际成本曲线的交点，交点的横坐标即最优计划量。
\textbf{两点限定务必保留}：其一，上述推导\textbf{不依赖任何分布假设}，
(a) 的经验分布仅示范形状、并非推导前提；其二，该结论\textbf{只在忽略储能与
跨时段耦合时严格成立}，故本文把它作为参数搜索的起点而非最终结论，最终由连续
运行总费用标定（8 轮滚动标定落在 $\alpha\in[0.70,0.90]$，其中 0.75 出现最多）。
经验分布由 334 天 $\times$ 144 时段共 48096 个样本汇总；模型内部实为\textbf{逐时段}
取 28 天窗口的 $\alpha$ 分位数。


% ---------------------------------------------------------------- 图 4
\begin{figure}[htbp]
  \centering
  \includegraphics[width=\textwidth]{fig04_risk_execution.png}
  \caption{问题二：实际净负荷超出风险修正净负荷，并不必然触发紧急购电}
  \label{fig:p2-trigger}
\end{figure}

\noindent 注：(a)(b) 为冬至 2025-12-21，(c)(d) 为秋分 2025-09-23；(a)(c) 为
净负荷，(b)(d) 为储电量。储备线 $R_t=E_{\min}+\rho(\bar E_t-E_{\min})$，
其中 $E_{\min}=1200$ kWh，$\rho=0.9$ 取自滚动标定结果；执行层放电不得使储电量
低于该线，故超出部分能否被吸收，取决于储能是否还有余量。
冬至当日有 \textbf{18} 个时段实际净负荷超出风险修正净负荷（最大超出 84.4 kWh），
但储能全程未触及下界（最低 1401 kWh），\textbf{紧急购电为 0}；
秋分当日有 \textbf{71} 个时段超出（最大超出 133.7 kWh），储能被放到 1200 kWh
下界，最终有 \textbf{19} 个时段发生紧急购电。
两天的"超出"都存在，唯一的差别是储能余量——这正是本文强调触发机制为
\textbf{"实际超出风险余量的部分"与"储能剩余可放能力"两者之差}，而非
"实际 $>$ 风险净负荷"本身的原因。
另需说明：秋分正午光伏出力超过负载，净负荷最低达 \textbf{$-479$ kWh}（26 个
时段为负），故图中纵轴完整保留了负值区间并绘出零线。


% ---------------------------------------------------------------- 图 5
\begin{figure}[htbp]
  \centering
  \includegraphics[width=\textwidth]{fig05_cost_waterfall.png}
  \caption{问题二：风险修正策略与预测均值策略的全年费用对照（瀑布图）}
  \label{fig:p2-waterfall}
\end{figure}

\noindent 注：两端柱为两策略各自的全年合计购电费，中间两根为二者之差。
与预测均值策略相比，风险修正策略的计划购电费\textbf{增加} 98.9 万元，
紧急购电费\textbf{减少} 195.3 万元，全年合计\textbf{净省 96.4 万元}。
关键在于它\textbf{并非买得更少}——计划购电量实际\textbf{增加 6.7\%}——
而是用低价的确定性支出，替代了 5 倍电价的尾部风险支出（紧急购电量
\textbf{下降 74.3\%}）。两策略使用同一预测器、同一实时执行规则，各自单独标定；
统计区间为 2025-02-01 至 2025-12-31，共 334 天。数据见支撑材料
\texttt{p2\_results.json}。
